本文围绕题目提出的主流大语言模型评价与选择问题展开。公开评测结果能够反映模型在不同任务上的表现，但各测试的量纲、覆盖范围和发布完整性并不一致；同一模型的公开成绩还可能因评测设置不同而出现缺失。因此，题目中的“综合评价”不能简单理解为把所有分数相加，而应先回答三个相互衔接的问题。

\subsection{问题一：主流模型的综合性能评价}

题目要求依据公开评测信息，对候选模型进行综合比较。其数学转译为：在模型--测试矩阵存在缺失、不同测试不可直接相加且同类测试可能重复计量的条件下，建立可比的多维能力评价体系，给出各模型的五维能力结构、综合得分和排序，并检验结论对样本与指标变化的敏感性。

问题一的输出包括复杂推理、知识与事实可靠性、长上下文、代码与软件工程、多模态五个能力维度的相对得分，以及基于正式权重的 Ranking A 综合结果。该结果回答的是“模型具有什么基础能力”，并不直接回答某一类任务中应当选择谁。

\subsection{问题二：不同任务场景下的模型选择}

综合能力排序只能给出总体评价，不能表达科研长文本分析、大众日常通用对话和计算机代码开发对能力组合的不同需求。题目进一步要求把模型能力与具体使用场景联系起来。其数学转译为：以问题一得到的五维能力为输入，根据题设场景需求确定能力权重，并用能够体现能力互补和短板约束的聚合函数计算场景效用。

问题二的输出是三类场景内的效用与排名。不同场景使用的能力集合和替代弹性不同，因而效用只在同一场景内比较。该结果回答的是“模型是否适合某类任务”，也解释了为什么综合排名靠前的模型在场景中可能发生名次迁移。

\subsection{问题三：成本约束下的部署决策}

即使模型在某一场景中具有较高效用，也不能据此直接作出部署选择。输入长度、输出长度、调用次数和 API 单价共同决定任务成本，且成本信息可能存在可核验性差异。题目要求在性能与成本之间进行权衡。其数学转译为：将问题二的场景效用与标准化工作负载成本联合起来，构造效用--成本 Pareto 前沿，并在不同预算上限下给出最优可行模型。

由此，三问构成递进关系：
\[
\text{公开评测信息}
\longrightarrow
\text{基础能力识别}
\longrightarrow
\text{场景适配}
\longrightarrow
\text{预算约束下的部署决策}.
\]
